\subsection{FILTERED BIGEBRAS}

DEFINITION 2. Let E be a bigebra with coproduct c.~A filtration compatible with the bigebra structure on E is an increasing sequence \((\mathrm{E}_{n})_{n\geq0}\) of submodules of E such that

\[
\begin{array}{c} \mathrm{E} _ {0} = \mathrm{K}. 1, \qquad \mathrm{E} = \bigcup_ {n \geqslant 0} \mathrm{E} _ {n} \\ \mathrm{E} _ {m}. \mathrm{E} _ {n} \subset \mathrm{E} _ {m + n} \quad \text {for} \quad m \geqslant 0, n \geqslant 0 \\ c (\mathrm{E} _ {n}) \subset \sum_ {i + j = n} \mathrm{Im} (\mathrm{E} _ {i} \otimes \mathrm{E} _ {j}) \quad \text {for} n \geqslant 0. \dagger \end{array}\tag{6}
\]

A bigebra with a filtration compatible with its bigebra structure is called a filtered bigebra.

Example. Let E be a graded bigebra (Algebra, Chapter III, § 11, no. 4, Definition 3) and \((\mathrm{E}^n)_{n\geqslant 0}\) its graduation. We write \(\mathrm{E}_n = \sum_{i = 0}^{n}\mathrm{E}^i\) . The sequence \((\mathrm{E}_n)\) is a filtration compatible with the bigebra structure on E.

PROPOSITION 6. Let E be a filtered bigebra and \((\mathrm{E}_n)_{n\geqslant 0}\) its filtration. For every integer \(n\geqslant 0\) , let \(\mathrm{E}_n^+ = \mathrm{E}_n\cap \mathrm{E}^+\) . Then \(\mathrm{E_0^+} = \{0\}\) and

\[
c ^ {+} (\mathrm{E} _ {n} ^ {+}) \subset \sum_ {i = 1} ^ {n - 1} \operatorname{Im} (\mathrm{E} _ {i} ^ {+} \otimes \mathrm{E} _ {n - i} ^ {+}) \quad \text { for } n \geqslant 0. \dagger\tag{7}
\]

As \(\mathrm{E}_0 = \mathrm{K}.1, \mathrm{E}_0^+ = 0\) . If \(x \in \mathrm{E}_n\) , \(\pi(x) = x - \varepsilon(x).1\) (formula (1)), whence \(\pi(x) \in \mathrm{E}_n^+\) and \(\pi(\mathrm{E}_n) \subset \mathrm{E}_n^+\) . It follows that \(\pi \otimes \pi\) maps \(\operatorname{Im}(\mathrm{E}_i \otimes \mathrm{E}_j)\) into \(\operatorname{Im}(\mathrm{E}_i^+ \otimes \mathrm{E}_j^+)\) for \(i \geqslant 0, j \geqslant 0\) . As \(c^+ = (\pi \otimes \pi) \circ c\) in \(\mathrm{E}^+\) (no. 1, Proposition 2), by (6)

\[
c ^ {+} \left(\mathrm{E} _ {n} ^ {+}\right) \subset \sum_ {i = 0} ^ {n} \operatorname{Im} \left(\mathrm{E} _ {i} ^ {+} \otimes \mathrm{E} _ {n - i} ^ {+}\right) = \sum_ {i = 1} ^ {n - 1} \operatorname{Im} \left(\mathrm{E} _ {i} ^ {+} \otimes \mathrm{E} _ {n - i} ^ {+}\right).
\]

COROLLARY. The elements of \(\mathbf{E}_1^+\) are primitive.

If \(x \in \mathrm{E}_{1}^{+}\) , then \(c^{+}(x) = 0\) by (7), whence (5).

